%==================================================================
% Ini adalah bab 2
% Silahkan edit sesuai kebutuhan, baik menambah atau mengurangi \section, \subsection
%==================================================================

\chapter[LANDASAN TEORI DAN KERANGKA PEMIKIRAN]{\\ LANDASAN TEORI DAN KERANGKA PEMIKIRAN}

\section{Tinjauan Pustaka}
\begin{sectioncontent}
    \hspace{\parindent}Penelitian mengenai klasifikasi penyakit daun tomat berbasis citra terus berkembang seiring dengan kemajuan \textit{computer vision} dan \textit{deep learning}. Berbagai penelitian terdahulu menunjukkan bahwa citra daun dapat digunakan sebagai dasar untuk mengenali kondisi sehat maupun gejala penyakit tanaman secara otomatis. Pendekatan yang digunakan juga semakin beragam, mulai dari arsitektur CNN umum, model ringan, hingga pendekatan yang diarahkan untuk penerapan pada perangkat dengan sumber daya terbatas.

    Kajian yang dilakukan oleh \citet{jackulin2022plantreview} menunjukkan bahwa \textit{machine learning} dan \textit{deep learning} telah banyak digunakan dalam deteksi penyakit tanaman melalui analisis citra daun. Pendekatan berbasis \textit{computer vision} dinilai mampu membantu proses identifikasi menjadi lebih konsisten dibandingkan pengamatan manual. Meskipun demikian, pembahasan pada kajian tersebut masih bersifat umum pada domain penyakit tanaman dan belum secara khusus membahas penyakit daun tomat maupun implementasi pada perangkat \textit{edge}.

    Penelitian yang lebih spesifik pada tanaman tomat dibahas oleh \citet{jelali2024tomatoreview}. Kajian tersebut menunjukkan bahwa \textit{deep learning} telah banyak dimanfaatkan untuk mendeteksi penyakit dan hama pada tanaman tomat, baik menggunakan dataset laboratorium maupun citra lapangan. Hal ini memperlihatkan bahwa tomat merupakan objek penelitian yang relevan dalam pengembangan model berbasis citra. Namun, fokus pembahasannya masih mencakup berbagai arsitektur secara umum dan belum secara khusus membandingkan dua generasi model ringan YOLO untuk tugas klasifikasi pada perangkat \textit{edge}.

    Ulasan lain dari \citet{das2025tomatoreview} memperlihatkan bahwa teknik klasifikasi, deteksi, dan segmentasi berbasis \textit{deep learning} telah memberikan kemajuan penting dalam identifikasi penyakit daun tomat. Kajian tersebut juga menekankan bahwa tantangan penelitian tidak hanya terletak pada pencapaian akurasi yang tinggi, tetapi juga pada kemampuan model untuk diterapkan dalam kondisi nyata. Kondisi lapangan umumnya memiliki keterbatasan perangkat, variasi lingkungan, serta kebutuhan respons yang cepat. Oleh karena itu, aspek efisiensi menjadi penting untuk diperhatikan selain akurasi model.

    Dari sisi ketersediaan data, beberapa dataset publik terbaru semakin memperkuat peluang penelitian pada domain penyakit daun tomat. Dataset yang disediakan oleh \citet{imtiaz2025dataset} memuat 731 citra resolusi tinggi daun tomat yang mencakup enam jenis penyakit umum dan daun sehat. Dataset dari \citet{oni2025dataset} memuat 1.028 citra daun tomat sehat dan bergejala yang diambil pada kondisi lapangan yang beragam. Sementara itu, \citet{liu2025dataset} menyediakan 1.026 citra penyakit tomat dari lingkungan rumah kaca modern dengan anotasi yang dikurasi untuk mendukung tugas klasifikasi maupun deteksi. Ketersediaan dataset tersebut menunjukkan bahwa penelitian klasifikasi penyakit daun tomat dapat dilakukan dengan data yang cukup jelas, terstruktur, dan terdokumentasi.

    Pada sisi pengembangan model, \citet{sun2024tomatoclassification} menunjukkan bahwa penggabungan strategi ekstraksi fitur dan pembelajaran mendalam dapat meningkatkan performa klasifikasi penyakit daun tomat. Akan tetapi, penelitian tersebut belum secara khusus membandingkan model keluarga YOLO varian \textit{nano}. Penelitian yang lebih dekat dengan konteks implementasi perangkat terbatas dilakukan oleh \citet{majeed2024edge}, yaitu dengan mengevaluasi solusi \textit{edge AI} mandiri untuk penyakit tomat menggunakan jaringan ringan. Penelitian tersebut memperkuat pandangan bahwa model untuk implementasi lapangan perlu dipilih berdasarkan keseimbangan antara akurasi dan efisiensi inferensi.

    Dalam konteks komputasi \textit{edge}, \citet{singh2023edge}, \citet{gill2025taxonomy}, dan \citet{xu2025edge} menjelaskan bahwa \textit{edge AI} menjadi pendekatan yang semakin penting karena memungkinkan proses inferensi dilakukan lebih dekat dengan sumber data. Pendekatan ini dapat mengurangi latensi, menurunkan ketergantungan terhadap komputasi awan, serta mendukung respons sistem yang lebih cepat. Dengan karakteristik tersebut, perangkat \textit{edge} menjadi relevan untuk sistem klasifikasi penyakit tanaman yang ditujukan untuk membantu proses pengambilan keputusan secara langsung di lapangan.

    Sementara itu, perkembangan keluarga YOLO dibahas oleh \citet{terven2023yolo}, \citet{ali2024yoloframework}, dan \citet{murat2025yolo}. Ketiga referensi tersebut menunjukkan bahwa YOLO berkembang sebagai keluarga model yang menekankan keseimbangan antara kecepatan dan performa. Pada awalnya, YOLO lebih dikenal dalam tugas \textit{object detection}. Namun, perkembangan ekosistem YOLO modern membuat model ini juga dapat digunakan untuk tugas lain, termasuk \textit{image classification}. Dokumentasi resmi Ultralytics menunjukkan bahwa YOLOv8 dan YOLO11 mendukung tugas klasifikasi citra, sehingga keduanya relevan untuk dibandingkan dalam penelitian ini \citep{ultralyticsclassify,ultralyticsyolov8,ultralyticsyolo11}.

    Berdasarkan kajian terhadap penelitian terdahulu, dapat disimpulkan bahwa klasifikasi penyakit daun tomat berbasis citra telah terbukti layak dilakukan menggunakan \textit{deep learning}. Selain itu, ketersediaan dataset publik mendukung pelaksanaan penelitian secara lebih terukur. Namun, aspek efisiensi implementasi pada perangkat \textit{edge} masih menjadi isu penting yang belum selalu dibahas secara mendalam. Penelitian yang secara khusus membandingkan YOLOv8 Nano dan YOLO11 Nano untuk klasifikasi penyakit daun tomat pada perangkat \textit{edge} juga masih relatif terbatas. Oleh karena itu, penelitian ini diarahkan untuk mengisi celah tersebut dengan membandingkan dua model ringan dari keluarga YOLO berdasarkan akurasi klasifikasi dan efisiensi komputasi.
\end{sectioncontent}

\section{Landasan Teori}
\begin{sectioncontent}
    \hspace{\parindent}Landasan teori digunakan untuk menjelaskan konsep-konsep utama yang menjadi dasar penelitian. Konsep tersebut meliputi penyakit daun tomat, klasifikasi citra, \textit{deep learning}, CNN, keluarga model YOLO, perangkat \textit{edge}, akuisisi citra, serta akurasi dan efisiensi sebagai dasar evaluasi. Pembahasan teori ini diperlukan agar metode penelitian dan analisis hasil memiliki dasar yang jelas.
\end{sectioncontent}

\subsection{Penyakit Daun Tomat}
\begin{subsectioncontent}
    \hspace{\parindent}Penyakit daun tomat merupakan gangguan pada daun tanaman tomat yang dapat disebabkan oleh jamur, bakteri, virus, maupun faktor lain. Gangguan tersebut umumnya memunculkan gejala visual tertentu, seperti bercak, perubahan warna, nekrosis, deformasi, atau kerusakan jaringan daun. Identifikasi dini terhadap gejala pada daun penting dilakukan karena daun sering menjadi indikator awal adanya gangguan kesehatan tanaman yang dapat memengaruhi pertumbuhan dan hasil panen \citep{jelali2024tomatoreview,das2025tomatoreview,nyawose2025plantreview}.

    Dalam penelitian berbasis citra, gejala visual pada daun tomat dapat dipelajari sebagai pola yang mewakili kelas penyakit tertentu. Dengan kata lain, citra daun tidak hanya dilihat sebagai gambar, tetapi juga sebagai data yang mengandung karakteristik visual untuk proses klasifikasi. Oleh karena itu, daun tomat menjadi objek yang sesuai untuk tugas klasifikasi citra, terutama apabila dataset yang digunakan telah memiliki label kelas yang jelas dan proses akuisisi citra dilakukan secara memadai \citep{imtiaz2025dataset,oni2025dataset,liu2025dataset}.
\end{subsectioncontent}

\subsection{Klasifikasi Citra}
\begin{subsectioncontent}
    \hspace{\parindent}Klasifikasi citra adalah proses pemberian label pada suatu gambar berdasarkan karakteristik visual yang terdapat di dalamnya. Pada tugas ini, citra diperlakukan sebagai satu kesatuan masukan, kemudian model menghasilkan satu label kelas sebagai keluaran. Dalam penelitian ini, label keluaran merepresentasikan jenis penyakit daun tomat atau kondisi daun sehat.

    Klasifikasi citra berbeda dengan \textit{object detection}. Pada \textit{object detection}, model tidak hanya menentukan jenis objek, tetapi juga mencari posisi objek tersebut di dalam citra. Sementara itu, klasifikasi citra hanya berfokus pada penentuan kategori gambar secara keseluruhan. Oleh karena itu, pendekatan klasifikasi dipilih dalam penelitian ini karena tujuan utamanya adalah menentukan kelas penyakit dari citra daun tomat, bukan menemukan lokasi area penyakit pada daun \citep{ultralyticsclassify,sun2024tomatoclassification}.
\end{subsectioncontent}

\subsection{\textit{Deep Learning} dan CNN}
\begin{subsectioncontent}
    \hspace{\parindent}\textit{Deep learning} merupakan cabang dari \textit{machine learning} yang menggunakan jaringan saraf tiruan berlapis untuk mempelajari pola kompleks dari data. Dalam pengolahan citra, pendekatan ini banyak digunakan karena mampu mengekstraksi fitur visual secara otomatis tanpa bergantung sepenuhnya pada rekayasa fitur manual \citep{jackulin2022plantreview,nyawose2025plantreview}.

    Salah satu arsitektur yang umum digunakan dalam klasifikasi citra adalah \textit{Convolutional Neural Network} (CNN). CNN memanfaatkan operasi konvolusi untuk menangkap pola spasial pada citra, seperti tepi, tekstur, warna, dan bentuk. Pola-pola tersebut kemudian diproses lebih lanjut untuk menghasilkan keputusan klasifikasi. Dalam penelitian penyakit tanaman, CNN banyak digunakan karena mampu mengenali gejala visual pada daun yang sering muncul dalam bentuk bercak, perubahan warna, atau pola tekstur tertentu \citep{jelali2024tomatoreview,das2025tomatoreview,sun2024tomatoclassification}.
\end{subsectioncontent}

\subsection{Keluarga Model YOLO}
\begin{subsectioncontent}
    \hspace{\parindent}YOLO (\textit{You Only Look Once}) merupakan keluarga model \textit{deep learning} yang awalnya dikenal luas pada tugas \textit{object detection}. Model ini banyak digunakan karena dirancang untuk menghasilkan prediksi secara cepat dalam satu proses inferensi. Perkembangan YOLO dari versi awal hingga versi yang lebih baru menunjukkan adanya peningkatan pada efisiensi, akurasi, dan fleksibilitas penggunaan \citep{terven2023yolo,ali2024yoloframework,murat2025yolo}.

    Pada perkembangannya, varian modern dari keluarga YOLO tidak hanya digunakan untuk \textit{object detection}, tetapi juga mendukung tugas lain, termasuk klasifikasi citra. Dokumentasi resmi Ultralytics menjelaskan bahwa \textit{image classification} digunakan untuk mengelompokkan keseluruhan citra ke dalam kelas yang telah ditentukan. Dukungan tersebut tersedia pada ekosistem YOLOv8 maupun YOLO11, sehingga kedua model dapat digunakan dalam penelitian klasifikasi citra \citep{ultralyticsclassify,ultralyticsyolov8,ultralyticsyolo11}.

    Dalam penelitian ini, model yang dibandingkan adalah YOLOv8 Nano dan YOLO11 Nano. Keduanya termasuk varian ringan atau \textit{nano} yang dirancang untuk memberikan keseimbangan antara ukuran model, kecepatan inferensi, dan performa prediksi. Karakteristik tersebut membuat keduanya relevan untuk diuji pada lingkungan komputasi terbatas seperti perangkat \textit{edge}. Selain itu, pemilihan YOLOv8 Nano dan YOLO11 Nano juga memungkinkan penelitian ini membandingkan dua generasi model ringan dari keluarga yang sama, sehingga analisis dapat lebih terarah pada perbedaan performa antarversi model.
\end{subsectioncontent}

\subsection{Perangkat \textit{Edge}}
\begin{subsectioncontent}
    \hspace{\parindent}Komputasi \textit{edge} adalah pendekatan pemrosesan data yang dilakukan lebih dekat dengan sumber data atau lokasi tempat data dihasilkan. Pada pendekatan ini, inferensi model dapat dijalankan secara lokal tanpa selalu bergantung pada server pusat atau komputasi awan. Keunggulan utama dari pendekatan ini adalah latensi yang lebih rendah, respons yang lebih cepat, serta potensi pengurangan komunikasi data ke server eksternal \citep{singh2023edge,gill2025taxonomy,xu2025edge}.

    Dalam penelitian ini, perangkat \textit{edge} digunakan sebagai media implementasi untuk menguji efisiensi model klasifikasi. Raspberry Pi 5 dipilih karena merupakan platform komputasi lokal yang mendukung CPU \textit{quad-core} Arm Cortex-A76 2.4 GHz dan varian RAM hingga 16 GB, termasuk versi 8 GB yang digunakan dalam penelitian ini \citep{raspberrypi5,raspberrypi5brief}. Pengujian pada perangkat \textit{edge} penting dilakukan karena model yang akurat belum tentu layak digunakan pada sistem nyata apabila waktu inferensinya terlalu tinggi atau penggunaan sumber dayanya terlalu besar.
\end{subsectioncontent}

\subsection{Akuisisi Citra dan Faktor Kamera}
\begin{subsectioncontent}
    \hspace{\parindent}Pada sistem klasifikasi berbasis citra, kualitas data masukan sangat memengaruhi hasil prediksi model. Faktor seperti resolusi gambar, pencahayaan, sudut pengambilan, jarak kamera terhadap objek, latar belakang, dan kestabilan posisi kamera dapat memengaruhi karakteristik visual citra. Pada citra daun tomat, variasi tersebut dapat memengaruhi keterbacaan gejala penyakit, seperti bercak, perubahan warna, dan tekstur permukaan daun \citep{imtiaz2025dataset,oni2025dataset,liu2025dataset}.

    Dalam penelitian ini, kamera tidak diposisikan sebagai variabel bebas, melainkan sebagai alat akuisisi citra yang kondisinya perlu dikendalikan. Artinya, proses pengambilan citra perlu dilakukan dalam kondisi yang dibuat konsisten agar pengaruh variasi kamera dan lingkungan dapat diminimalkan. Dengan pengendalian tersebut, perbandingan model dapat dilakukan secara lebih objektif karena perbedaan hasil pengujian lebih mencerminkan performa model daripada perbedaan kualitas data akuisisi.
\end{subsectioncontent}

\subsection{Akurasi dan Efisiensi sebagai Dasar Evaluasi}
\begin{subsectioncontent}
    \hspace{\parindent}Akurasi dalam penelitian klasifikasi menunjukkan tingkat ketepatan model dalam memberikan label yang benar terhadap data uji. Selain akurasi, metrik lain seperti \textit{precision}, \textit{recall}, dan \textit{F1-score} juga digunakan untuk memberikan gambaran performa yang lebih rinci. Pada klasifikasi multikelas, metrik tersebut penting karena model tidak hanya perlu memiliki akurasi keseluruhan yang baik, tetapi juga perlu mampu membedakan setiap kelas penyakit secara konsisten.

    Efisiensi menjadi aspek penting ketika model diterapkan pada perangkat \textit{edge}. Efisiensi dapat dinilai melalui beberapa indikator, seperti waktu inferensi, penggunaan memori, ukuran model, dan penggunaan sumber daya komputasi lain yang relevan. Dalam penelitian ini, akurasi dan efisiensi digunakan sebagai dua dasar utama untuk menentukan model yang lebih sesuai diterapkan pada sistem klasifikasi penyakit daun tomat berbasis perangkat \textit{edge} \citep{singh2023edge,gill2025taxonomy,xu2025edge,majeed2024edge}.
\end{subsectioncontent}

\newpage

\section{Kerangka Pemikiran}
\begin{sectioncontent}
    % \hspace{\parindent}Kerangka pemikiran dalam penelitian ini disusun untuk menggambarkan alur hubungan antara permasalahan, dasar teori, model yang digunakan, proses pengujian, hingga penentuan model yang lebih sesuai. Penelitian ini berangkat dari permasalahan penyakit daun tomat yang membutuhkan proses identifikasi yang cepat dan konsisten. Pendekatan berbasis citra dan \textit{deep learning} dipilih karena mampu mempelajari pola visual pada daun, sedangkan perangkat \textit{edge} digunakan untuk menilai kelayakan model pada lingkungan komputasi terbatas.

    % Berdasarkan kajian pustaka dan landasan teori, YOLOv8 Nano dan YOLO11 Nano dipilih sebagai dua model ringan yang akan dibandingkan. Kedua model dilatih dan diuji menggunakan dataset yang sama, kemudian hasilnya dianalisis berdasarkan akurasi klasifikasi dan efisiensi inferensi. Alur kerangka pemikiran dalam penelitian ini ditunjukkan pada Gambar \ref{fig:kerangka-pemikiran}.

    \begin{figure}[H]
        \centering
        \includegraphics[width=0.8\textwidth,height=0.8\textheight]{kerangka-pemikiran.png}
        \caption{Kerangka Pemikiran}
        \label{fig:kerangka-pemikiran}
    \end{figure}
\end{sectioncontent}